\documentclass[UTF8,a4paper,10pt,fontset=none]{ctexart}

\usepackage{kexiearticle}

\title{How Far Can Heuristic Learning Go? 

Characterizing Performance Ceilings,
Complexity, and Sample Efficiency in Large-Scale Real-World Competitive Games}
\shorttitle{How Far Can Heuristic Learning Go?}
\author{Kaisen Yang$^{\sharp}$}
% @刘青乐确认一下author list
\affiliation{Department of Computer Science and Technology, Tsinghua University}
\role[\sharp]{Corresponding Author.}
\email{yks23@mails.tsinghua.edu.cn}
\articledate{2026-08-15}

\begin{document}

\makekexietitle

\begin{abstract}
% Write the abstract here.
\end{abstract}
% 需要一个头图，大概表示我们干了什么事情
\section{Introduction}
% contribution: 1. 定义了HL范式及其优化过程 2. 建立了benchmark 3. 验证了xxx等现象
\section{Preliminary}
% 定义清楚要用到的东西
\section{Method}
% method图，展示算法核心概念
\subsection{Motivation}
% 何以想到HL的？- Agent - 样本Efficiency - 实际的效果
\subsection{HL Algorithms}
% 定义清楚有哪些自己的HL算法
\section{Experiments}
\subsection{set up of Benchmark}
% Benchmark的示意图
% 先把讨论的实验做了搬上来
\subsection{Experiments}

\section{Related Work}
\section{Future work and Limitation}
\section{Conclusion}

\appendix

\end{document}
